% Chapter 8: Conclusion

\chapter{Conclusion}
\label{ch:conclusion}

This thesis asked how controlled changes to visual evidence affect answers and
rationales in a two-stage VQA pipeline. We evaluated four separately adapted
input arms on the same 2,653 VCR records under source, grey, mismatch, mask,
noise, and mirror conditions.

RQ1 showed that Stage~1 answer behaviour depended on the available visual
evidence. Mismatch caused the largest accuracy loss, followed by grey,
then mask and noise; mirror left net accuracy close to source while still
causing correctness changes for individual records. Description arms retained
more predictions that were correct under source when images were degraded, a
pattern consistent with the fixed description remaining available as a text
channel.

RQ2 showed that Stage~2 rationale text conditioned on the gold answer was
sensitive to the same interventions. Change was measured with BERTScore on the
final rationale span, and followed the same broad order; the reasoning span and
the complete generation give the same ordering
(Appendix~\ref{sec:app-rq2-sensitivity}).
The matched control with a wrong answer produced greater Drift than
matched grey removal in every arm, showing that rationale generation was also conditioned by
the answer it was asked to justify.

A secondary predicted-answer regime also showed that coverage of usable
rationales can vary sharply between arms, and that recovering the supplied answer from a
rationale does not establish that the answer is correct.

The most consistent robustness pattern across the two stages was associated with
descriptions. Polygon overlays alone did not yield a consistent independent
advantage.

These results are behavioural evidence relevant to faithfulness. They show how
outputs respond to controlled changes in visual evidence and supplied answers,
within the limits set out in Sections~\ref{sec:discussion-faithfulness}
and~\ref{sec:discussion-limitations}.

Arms with a fixed textual description appeared more robust under image
degradation at both stages, a pattern consistent with scene information remaining
available through text. The arms that held up best under grey and mismatch had an
additional route through text. Their robustness therefore does not establish
stronger visual grounding. An evaluation that varies only the image and reports
only accuracy will therefore reward auxiliary text channels and cannot distinguish
them from visual evidence use. Separating the
two requires varying the text channel as well. This applies to the system
evaluated in this thesis; how far it carries to other models and datasets is an
open question.

Future work should repeat the six image conditions across models, datasets, and
training seeds. Independent human validation would strengthen the interpretation
of metrics for rationale change. A design that varies description availability
systematically during training and evaluation could separate what a checkpoint
expects from what it substitutes at inference. Internal evidence use will
require causal or mechanistic methods beyond output comparison.
